%% ringraziamenti.tex
%%
%% Copyright (C) 2000-2018 Simone Piccardi.  Permission is granted to
%% copy, distribute and/or modify this document under the terms of the GNU Free
%% Documentation License, Version 1.1 or any later version published by the
%% Free Software Foundation; with the Invariant Sections being "Un preambolo",
%% with no Front-Cover Texts, and with no Back-Cover Texts.  A copy of the
%% license is included in the section entitled "GNU Free Documentation
%% License".
%%
\chapter{Ringraziamenti}
\label{cha:ringraziamenti}

Desidero ringraziare tutti coloro che a vario titolo e a più riprese mi hanno
aiutato ed han contribuito a migliorare in molteplici aspetti la qualità di
GaPiL. In ordine rigorosamente alfabetico desidero citare:
\begin{description}
\item[\textbf{Alessio Frusciante}] per l'apprezzamento, le innumerevoli
  correzioni ed i suggerimenti per rendere più chiara l'esposizione.
\item[\textbf{Daniele Masini}] per la rilettura puntuale, le innumerevoli
  correzioni, i consigli sull'esposizione ed i contributi relativi alle
  calling convention dei linguaggi e al confronto delle diverse tecniche di
  gestione della memoria.
\item[\textbf{Mirko Maischberger}] per la rilettura, le numerose correzioni,
  la segnalazione dei passi poco chiari e soprattutto per il grande lavoro
  svolto per produrre una versione della guida in un HTML piacevole ed
  accurato.
\item[\textbf{Fabio Rossi}] per la rilettura, le innumerevoli correzioni, ed i
  vari consigli stilistici ed i suggerimenti per il miglioramento della
  comprensione di vari passaggi.
\end{description}

Infine, vorrei ringraziare il Firenze Linux User Group (FLUG), di cui mi
pregio di fare parte, che ha messo a disposizione il repository CVS su cui era
presente la prima versione della Guida, ed il relativo spazio web, e
\href{https://www.truelite.it}{Truelite Srl}, l'azienda che ho fondato e di
cui sono amministratore, che fornisce le risorse per il nuovo repository Git
ed per il sito del progetto: \url{https://gapil.gnulinux.it}.


% LocalWords:  GaPiL Masini calling convention Maischberger HTML Group FLUG CVS
% LocalWords:  repository Truelite Srl SVN


%%% Local Variables: 
%%% mode: latex
%%% TeX-master: "gapil"
%%% End: 

